% Archivo compartido: no compilar directamente.
\ifdefined\documentoSemanal
\else
\usepackage{tikz}
\usepackage{estilo_apuntes}
\usetikzlibrary{babel}
\encabezado{Semana 2 · Clase 3}{Material de clase}
\begin{document}
\fi

\ifdefined\documentoSemanal
\else
\tituloclase{Límites básicos y álgebra de límites}{Sustitución, factorización y racionalización}{Semana 2 · Clase 3 · 2 horas}

\begin{objetivos}
\begin{itemize}
  \item Aplicar las leyes algebraicas de los límites.
  \item Calcular límites de funciones polinómicas, racionales y radicales.
  \item Reconocer cuándo la sustitución directa es válida.
  \item Resolver indeterminaciones $0/0$ mediante factorización o racionalización.
\end{itemize}
\end{objetivos}

\begin{resumen}
Las leyes de límites permiten trasladar operaciones algebraicas de las funciones a sus límites. La sustitución directa funciona para las funciones elementales en puntos de su dominio. Si produce $0/0$, esa expresión no es la respuesta: indica que debe buscarse una función localmente igual mediante simplificación.
\end{resumen}
\fi

\section*{Límites básicos}

\begin{proposicion}[límites fundamentales]
Sean $a,c\in\mathbb{R}$ y $n\in\mathbb{N}^{*}$. Entonces
\[
\lim_{x\to a}c=c,\qquad
\lim_{x\to a}x=a,\qquad
\lim_{x\to a}x^n=a^n.
\]
Cuando la raíz está definida en un entorno adecuado de $a$,
\[
\lim_{x\to a}\sqrt[n]{x}=\sqrt[n]{a}.
\]
Para $n$ par y $a=0$, esta última afirmación se interpreta como límite por la derecha.
\end{proposicion}

\begin{teorema}[álgebra de límites]
Suponga que $\lim_{x\to a}f(x)=L$, $\lim_{x\to a}g(x)=M$ y $\lambda\in\mathbb{R}$. Entonces:
\begin{align*}
\lim_{x\to a}(f(x)+g(x))&=L+M,\\
\lim_{x\to a}(f(x)-g(x))&=L-M,\\
\lim_{x\to a}\lambda f(x)&=\lambda L,\\
\lim_{x\to a}f(x)g(x)&=LM.
\end{align*}
Si $M\neq0$, también
\[
\lim_{x\to a}\frac{f(x)}{g(x)}=\frac{L}{M}.
\]
Además, las potencias enteras positivas y las raíces compatibles con su dominio pueden trasladarse al límite.
\end{teorema}

\begin{nota}
La condición $M\neq0$ en la regla del cociente es indispensable. Una forma $0/0$ es \textbf{indeterminada}: distintos cocientes pueden presentar esa forma y tener límites diferentes o no tener límite.
\end{nota}

\section*{Sustitución directa}

\begin{proposicion}[funciones algebraicas]
Sea $a\in\mathbb{R}$.
\begin{enumerate}
  \item Si $p$ es un polinomio, entonces $\lim_{x\to a}p(x)=p(a)$.
  \item Si $r=p/q$ es racional y $q(a)\neq0$, entonces $\lim_{x\to a}r(x)=r(a)$.
  \item Para expresiones radicales, la sustitución es válida cuando $a$ y los valores próximos considerados pertenecen al dominio de la raíz.
\end{enumerate}
\end{proposicion}

\begin{ejemplo}[sustitución directa]
Calcule
\[
\lim_{x\to-2}(3x^3-x+4),\qquad
\lim_{x\to1}\frac{x^2+2}{x+3},\qquad
\lim_{x\to5}\sqrt{x+4}.
\]
\end{ejemplo}
\ifmostrarSoluciones
\begin{solucion}
Los denominadores y radicales son admisibles en los puntos indicados. Por sustitución,
\[
-18,\qquad \frac34,\qquad 3,
\]
respectivamente.
\end{solucion}
\else\vspace{15mm}\fi

\section*{Indeterminaciones algebraicas}

\begin{proposicion}[estrategia de simplificación local]
Si la sustitución produce $0/0$:
\begin{enumerate}
  \item factorice numerador y denominador o racionalice la expresión;
  \item simplifique solamente para $x\neq a$;
  \item calcule el límite de la expresión simplificada;
  \item conserve las restricciones del problema original.
\end{enumerate}
La justificación es que ambas expresiones son localmente iguales alrededor de $a$.
\end{proposicion}

\begin{ejemplo}[factorización]
Calcule
\[
\lim_{x\to3}\frac{x^2-9}{x-3}.
\]
\end{ejemplo}
\ifmostrarSoluciones
\begin{solucion}
La sustitución produce $0/0$. Para $x\neq3$,
\[
\frac{x^2-9}{x-3}=\frac{(x-3)(x+3)}{x-3}=x+3.
\]
Las funciones son localmente iguales alrededor de $3$; por tanto, el límite es $6$.
\end{solucion}
\else\vspace{16mm}\fi

\begin{ejemplo}[racionalización del numerador]
Calcule
\[
\lim_{x\to0}\frac{\sqrt{x+4}-2}{x}.
\]
\end{ejemplo}
\ifmostrarSoluciones
\begin{solucion}
Se multiplica por el conjugado:
\begin{align*}
\frac{\sqrt{x+4}-2}{x}
&=\frac{(\sqrt{x+4}-2)(\sqrt{x+4}+2)}{x(\sqrt{x+4}+2)}\\
&=\frac{1}{\sqrt{x+4}+2},\qquad x\neq0.
\end{align*}
En consecuencia, el límite es $1/4$.
\end{solucion}
\else\vspace{18mm}\fi

\begin{ejemplo}[racionalización del denominador]
Calcule
\[
\lim_{x\to4}\frac{x-4}{\sqrt{x}-2}.
\]
\end{ejemplo}
\ifmostrarSoluciones
\begin{solucion}
Para $x\neq4$,
\[
\frac{x-4}{\sqrt{x}-2}
=\frac{(\sqrt{x}-2)(\sqrt{x}+2)}{\sqrt{x}-2}
=\sqrt{x}+2.
\]
Luego el límite es $4$.
\end{solucion}
\else\vspace{15mm}\fi

\begin{nota}
No debe cancelarse un sumando. Solo pueden cancelarse factores no nulos. Por ejemplo, en $(x^2-9)/(x-3)$ se factoriza primero y se cancela el factor $x-3$ bajo la condición $x\neq3$.
\end{nota}

\begin{ejercicios}[álgebra y cálculo de límites]
Calcule los siguientes límites e indique la técnica principal.
\begin{enumerate}
  \item $\displaystyle\lim_{x\to2}(4x^3-3x+1)$.
  \item $\displaystyle\lim_{x\to-1}\frac{x^2+3x+5}{x^2+4}$.
  \item $\displaystyle\lim_{x\to4}\sqrt{2x+1}$.
  \item $\displaystyle\lim_{x\to2}\frac{x^2-4}{x-2}$.
  \item $\displaystyle\lim_{x\to-3}\frac{x^2+5x+6}{x+3}$.
  \item $\displaystyle\lim_{x\to1}\frac{x^3-1}{x-1}$.
  \item $\displaystyle\lim_{x\to0}\frac{\sqrt{1+x}-1}{x}$.
  \item $\displaystyle\lim_{x\to5}\frac{x-5}{\sqrt{x-1}-2}$.
  \item $\displaystyle\lim_{x\to2}\frac{\sqrt{x+2}-2}{x-2}$.
  \item $\displaystyle\lim_{x\to0^+}\sqrt{x}$.
\end{enumerate}
\end{ejercicios}

\ifmostrarSoluciones
\begin{solucion}[de los ejercicios]
\begin{enumerate}
  \item $27$, por sustitución directa.
  \item $3/5$, por sustitución directa.
  \item $3$, por sustitución directa.
  \item $4$, pues $(x^2-4)/(x-2)=x+2$ para $x\neq2$.
  \item $-1$, pues $(x+2)(x+3)/(x+3)=x+2$ para $x\neq-3$.
  \item $3$, usando $x^3-1=(x-1)(x^2+x+1)$.
  \item $1/2$, al racionalizar el numerador.
  \item $4$, porque el cociente se simplifica a $\sqrt{x-1}+2$.
  \item $1/4$, al racionalizar el numerador.
  \item $0$; solo se considera la aproximación por la derecha.
\end{enumerate}
\end{solucion}
\fi

\ifdefined\documentoSemanal
\else
\fuentes{\textbf{Para ampliar:} \emph{Cálculo en una variable: notas de clase 2018-B}, apartado de álgebra de límites; J. Stewart, \emph{Cálculo de una variable: trascendentes tempranas}, §§2.2--2.3; y G. Strang--E. Herman, \emph{Cálculo, volumen 1}, OpenStax, §§2.2--2.3.}
\end{document}
\fi
